\section{Výkonnostní měření synchronizace}
\label{sec:evaluace:benchmarks}

Měření výkonu synchronizace dat pro repozitář \texttt{rust-lang/rust} probíhalo na lokální pracovní stanici. Pro zajištění reprodukovatelnosti experimentů je níže uvedena hardwarová a~síťová specifikace testovacího prostředí:

\begin{itemize}
  \item \textbf{Procesor:} AMD Ryzen 5 5600G (6 jader, 12 vláken, základní frekvence 3,9\,GHz)
  \item \textbf{Operační paměť:} 32\,GB DDR4 (3\,600\,MHz)
  \item \textbf{Úložiště:} 512\,GB NVMe SSD (KXG50ZNV512G, 3200 MBps čtení / 1050 MBps zápis)
  \item \textbf{Připojení k~síti:} Symetrické internetové připojení s~propustností 300\,Mb/s
\end{itemize}

Časová náročnost synchronizace je úzce spjata s~propustností disku a~parametry sítě. Získané výsledky proto slouží primárně k~relativnímu porovnání v~rámci této konkrétní konfigurace a~absolutní hodnoty se mohou při použití odlišného hardwaru měnit.

\subsection{Počáteční inicializace dat}
\label{sec:evaluace:benchmarks:inicializace}

Počáteční synchronizace slouží k prvotnímu stažení všech dat z repozitáře za zvolené časové období.
Měření bylo provedeno pro data od 1.\,1.\,2023, tedy za přibližně 3 roky a 4,5 měsíce.
Celková doba počáteční synchronizace činila přibližně 49 minut.
Podrobný přehled jednotlivých fází je uveden v tabulce~\ref{tab:evaluace:benchmarks:inicializace}.

\begin{table}[h]
  \centering
  \caption{Časy jednotlivých fází počáteční synchronizace (rust-lang/rust, od 2023-01-01)}
  \label{tab:evaluace:benchmarks:inicializace}
  \begin{tabular}{lr}
    \toprule
    \textbf{Fáze} & \textbf{Trvání} \\
    \midrule
    Stažení \PR z~GitHub \API & 15 min 35 s \\
    Vložení přispěvatelů z~\PR do databáze & $<$ 1 s \\
    Zpracování \PR a~výpočet aktivity souborů & 15 min 59 s \\
    Hromadné vložení \PR do databáze & 2 s \\
    Hromadné vložení aktivity souborů do databáze & 15 s \\
    Stažení issues z~GitHub \API & 17 min 17 s \\
    Vložení přispěvatelů z~issues do databáze & $<$ 1 s \\
    Vložení issues do databáze & 1 s \\
    Synchronizace týmů & $<$ 1 s \\
    \midrule
    \textbf{Celkem} & $\approx$ \textbf{49 min} \\
    \bottomrule
  \end{tabular}
\end{table}

Za uvedené období bylo staženo přibližně 64\,000 \PR a issues, k čemuž bylo zapotřebí přibližně
640 \API požadavků na GitHub (při maximální velikosti stránky 100 záznamů).
Synchronizace dat týmů vyžadovala jediný požadavek na \API rust týmu \sectionref{sec:rust_team_api}, který nepodléhá běžnému limitu.

% Počáteční synchronizaci je doporučeno spouštět jako samostatný proces -- nejlépe s vyhrazeným přístupovým tokenem -- a to pouze jednou při prvním nasazení.
% Alternativně lze zvolit kratší časové okno (například posledních 30 dní) a teprve poté spustit stahování kompletní historie událostí jako oddělený proces na pozadí, čímž se výrazně zkrátí doba inicializace.

Stahování \PR a~issues by bylo teoreticky možné urychlit paralelizací požadavků. Vzhledem k~limitům počtu požadavků \sectionref{subsec:implementation:gotchas:rate_limits} a~riziku dočasného zablokování přístupového tokenu bylo zvoleno sekvenční zpracování. Počáteční synchronizace se navíc v~ideálním případě provede pouze jednou při prvním nasazení, takže tato volba nemá zásadní vliv na běžný provoz systému.

\subsection{Inkrementální aktualizace}
\label{sec:evaluace:benchmarks:inkrementalni}

Inkrementální aktualizace je realizována pomocí timestampu poslední již uložené události v databázi. Systém si při zahájení synchronizace načte nejnovější známý timestamp a následně volá výpisy pull requestů a issues seřazené podle času poslední změny v sestupném pořadí z GitHub \API. Zpracovávány jsou pouze ty stránky výsledků, jejichž položky mohou ještě obsahovat záznamy novější než uložené mez.
Při dosažení starších položek se procházení ukončí.
Detailní historie událostí a změn štítků je potom stahována už jen pro takto vybrané změněné pull requesty a issues, přičemž i zde lze načítání ukončit po dosažení již dříve uložených událostí.

Po dokončení počáteční synchronizace systém provádí inkrementální aktualizace, při nichž jsou z GitHub \API staženy pouze záznamy změněné od posledního běhu.
Výpis~\ref{lst:evaluace:benchmarks:inkrementalni} zachycuje log inkrementální aktualizace v situaci, kdy nebyly nalezeny žádné nové záznamy.

\begin{listing}[h]
\begin{minted}[linenos=false]{text}
Getting pull requests from github:  took: 3.653 seconds
Inserting contributors from pull requests:  took: 0.000 seconds
Inserting pull requests to database:  took: 0.000 seconds
Getting issues from github:  took: 2.473 seconds
Inserting contributors from issues to database:  took: 0.000 seconds
Inserting issues to database:  took: 0.000 seconds
Upserting users from rust teams:  took: 0.069 seconds
Overall getting resources:  took: 6.312 seconds
\end{minted}
  \caption{Log inkrementální aktualizace bez nových záznamů}
  \label{lst:evaluace:benchmarks:inkrementalni}
\end{listing}

Celý cyklus inkrementální aktualizace trval přibližně 6,3 sekundy.
Největší část vykazovalo dotazování GitHub \API, konkrétně přibližně 3,7 sekundy pro \PR a 2,5 sekundy pro issues.
V případě nových záznamů přibývá průměrně 1,5 dodatečného \API požadavku na každý nový \PR nebo issue, a to z důvodu stahování kompletní historie událostí a štítků.

Po dokončení každého synchronizačního cyklu jsou aktualizovaná data okamžitě dostupná v dashboardu.
Maximální prodleva mezi vznikem změny v repozitáři a jejím zobrazením odpovídá nakonfigurovanému intervalu mezi synchronizacemi.
Při výchozím nastavení intervalu na 2 minuty je tato prodleva v nejhorším případě 2 minuty.

\subsection{Spotřeba API požadavků}
\label{sec:evaluace:benchmarks:api}

GitHub omezuje počet \API požadavků na 5\,000 za hodinu pro autentizované požadavky~\cite{github_rate_limits} (viz \ref{subsec:github_api_limits}).
Při standardním provozu spotřebuje systém minimálně 60 požadavků za hodinu, přičemž se vychází z přibližně 2 požadavků na jeden synchronizační cyklus (jeden pro \PR, jeden pro issues) při intervalu 2 minuty (30 cyklů za hodinu).
K tomu je nutné připočítat průměrně 1,5 požadavku za každý nový \PR nebo issue potřebný k synchronizaci historie.
Pro repozitář \texttt{rust-lang/rust} většinou přibývá v závislosti na denní době 0 až 5 nových záznamů za hodinu.

Celková průměrná spotřeba při normálním provozu za hodinu je tedy:
\[
  30 \times 2 + 5 \times 1{,}5 = 67{,}5 \text{ požadavků/hod}
\]

Tato hodnota představuje přibližně $1{,}4$~\% stanoveného limitu, takže systém může bez problémů fungovat i s výrazně kratšími synchronizačními intervaly nebo může být provozováno více instancí souběžně pro různé repozitáře se stejným přístupovým tokenem.
